\sectionguard
\section*{The Propers: Interpretive Possibilities}

\begin{studybox}{Exploratory editorial and AI proposals}
Everything below is exploratory editorial and AI proposal. None of it is
sourced historical intent, attributed teaching, or documented reception, and
none of it should be quoted as any of those. Each proposal joins at least two
precisely named appointed elements and ends on the strongest thing that stands
against it. \textbf{No connection here is claimed to be unknown, unprecedented,
first, or authored by the model}; where a proposal and a checked witness
conflict, the witness wins. The precedent search behind each of these, its
corpus and its bounds, is recorded in \texttt{research/scope.md} and summarised
in the terminal appendix.
\end{studybox}

\begin{proposal}{P1. Goods are ordered, not refused}
\pfield{Anchors}{Introit \cue{Int.}, \latin{mélior est dies una in átriis tuis
super mília}; Gradual \cue{Grad.}, \latin{bonum est confidere in Dómino, quam
confidere in hómine}; Gospel \cue{Gosp.}, \latin{nemo potest duóbus dóminis
servíre} and \latin{non potéstis Deo servíre et mammónæ}; Communion
\cue{Comm.}, \latin{Primum}; Collect \cue{Coll.}}
\pfield{Mechanism}{The formulary states one choice in two grammatical
registers. Seven comparatives of degree in four Latin constructions stand in
exactly three of the ten elements --- the Introit's \latin{mélior \dots\ super
mília}, the Gradual's \latin{Bonum est \dots\ quam} twice, the Gospel's
\latin{plus quam} twice, its \latin{Nonne \dots\ magis pluris} and its
\latin{quanto magis} --- while the Epistle, the three orations, the Alleluia,
the Offertory and the Communion carry none. The readings state the same choice
\textbf{exclusively}: \latin{hæc enim sibi ínvicem adversántur}, no man can
serve two masters, you cannot serve God and mammon. \textbf{The Gospel sits on
both sides of the split}, its \latin{Ideo} making the exclusive premise the
ground of the comparative argument, and the comparative register is what keeps
the exclusive one from reading as contempt for created goods.}
\pfield{Fruit}{A doctrine of ordered preference preached without contempt for
creation, and the reason this Mass can command an absolute exclusion and still
promise that all things shall be added.}
\pfield{What the element-by-element reading misses}{Read one element at a time,
the chants look like decoration around two hard sayings. Read together, the two
registers are doing different grammatical work on the same choice.}
\pfield{Strongest limit}{\textbf{The defeater is not a modern philological
worry: a named medieval expositor stands on both sides of it.} The \latin{quam}
of Ps.\ 117:8--9 may render a Hebrew idiom of \emph{exclusion} rather than of
degree, and \latin{super mília} is a Hebraism too, so three of the seven
comparatives may not be comparatives at all --- only the Gospel's four are
Greek-derived and immune. \textbf{Bruno the Carthusian asks exactly that
question about exactly these words, and answers it twice}: first that the word
is not set down as a comparative, so that trust in the Lord is useful and
recourse to the persecutor is not; then that it is comparative, \textbf{but only
if the princes are heavenly spirits and spiritual men}, on the ground that God
can help without them and they cannot help without God. His wording is not
quoted here, at either answer, the passage standing in a column that must be
read at a page image before any of it is quoted. \textbf{And a Greek Father at
the appointed verses denies the comparison outright.} John Chrysostom, on
Ps.\ 117:8--9 at \work{Patrologia Graeca} 55:331, holds that Scripture uses the
comparative form of things that admit no comparison, out of condescension to
weak hearers --- \emph{ouk ara synkrinon touto phesin, alla synkatabainon
ekeinois} --- so that on his reading the Gradual's twice-repeated \latin{Bonum
est \dots\ quam} is not a weighing of two goods at all. \textbf{So this
proposal is stated as contingent on the comparative reading of the Latin
\latin{quam}, and on Bruno's second answer rather than his first, and it must
answer Chrysostom rather than gesture at him; nothing in a Latin distribution
can settle which is right.} Secondly, the
lexical evidence reaches only part of the anchor set --- \textbf{the Communion
and the Collect are thematic members here and not lexical ones}, and that is
said rather than blurred. Thirdly, this must not be built as ``the chants soften
the readings'': it is one unit, not a two-sided contrast.}
\end{proposal}

\begin{proposal}{P2. What man cannot add}
\pfield{Anchors}{Gospel \cue{Gosp.}, Mt.\ 6:27 \latin{Quis autem vestrum
cógitans potest adícere ad statúram suam cúbitum unum?} and Mt.\ 6:33
\latin{adiciéntur}; Communion \cue{Comm.}, \latin{adiciéntur}; Collect
\cue{Coll.}, \latin{quia sine te lábitur humána mortálitas}; Epistle
\cue{Ep.}, \latin{ut non quæcúmque vultis, illa faciátis}}
\pfield{Mechanism}{One verb carries the exchange. The Gospel denies that any
man can \emph{add} a cubit to his stature and then promises that all things
\emph{shall be added}; the Communion repeats the promise alone; and the Collect
and the Epistle supply the same incapacity in other words. \textbf{What man
cannot add to himself, God adds to him}, and the reversal is in the morphology
--- active of the man who cannot, passive of the God who does.}
\pfield{Fruit}{A doctrine of grace read off a single verb, and the reason the
Mass's petition and its promise are one movement rather than two.}
\pfield{What the element-by-element reading misses}{Both occurrences of the
verb sit inside one pericope, so an element-by-element reading files them as a
feature of Matthew rather than as the axis on which the Collect's diagnosis and
the Communion's promise turn.}
\pfield{Strongest limit}{\textbf{More of this is documented reception than
proposal, and the proposal is what is left when that is subtracted.} Aquinas
already reads the cubit as an argument from providence at the same verse;
pseudo-Bede reads the Communion's own verb, \latin{Non dixit nobis dari, sed
adiici}; and Christian of Stavelot supplies its elided object with its limit,
\latin{Ad necessitatem, non ad superfluitatem}. \textbf{What remains
exploratory is only the cross-element join} --- the incapacity at v.~27 to the
promise six verses later, to the Communion's repetition of the promise alone,
and to the Collect's \latin{sine te lábitur humána mortálitas}. \textbf{Neither
the Collect nor the Epistle uses the verb}, so two of the four members join
doctrinally and not lexically, \textbf{a weaker join than P4's and stated as
such.} The Benziger \latin{editio iuxta typicam} spells the root with a doubled
\latin{-ii-} in all three places, and the difference does not touch the
argument. Its form at Mt.\ 6:27 is \latin{adiícere}, read on the Internet
Archive item's own single-page JP2 derivative and so \textbf{a reading of an
unregistered derivative of a registered item}; the doubled \latin{-ii-} at
Mt.\ 6:33 and at the Communion stands in that item's uncorrected optical text
layer, \textbf{no page of the registered page-image artifact having been
opened.} \textbf{And Chrysostom rests the promise on the same verb in the
opposite direction from this reading}, in Greek at \work{Patrologia Graeca}
57:303: he said not \emph{shall be given} but \emph{shall be added}, ``that you
may learn that the things present are nothing great, set beside the greatness of
the things to come,'' in the \work{Nicene and Post-Nicene Fathers} English,
first series --- for him the point of
\emph{added} is the smallness of what is added and not the passivity of the
recipient. \textbf{Finally, the passive with a divine subject is the settled
grammar of the Roman collect}, so the orations' half of this may be a fact about
how Latin orations are written; the counter-test would need a concordance over
the Missal's orations, and the optical layers available recover no multi-word
phrase of this formulary at all.}
\end{proposal}

\begin{proposal}{P3. Propitiation and salvation saturate the prayers the people do not sing}
\pfield{Anchors}{Collect \cue{Coll.}, \latin{propitiatióne perpétua};
Secret \cue{Sec.}, \latin{tuæ propitiátio potestátis}; Postcommunion
\cue{Postcomm.}, \latin{perpétuæ \dots\ salvatiónis efféctum}; Alleluia
\cue{All.}, \latin{Deo salutári nostro}}
\pfield{Mechanism}{The three orations and the Alleluia carry vocabulary no
other element does --- \latin{propitiátio} twice, the \latin{salut-} and
\latin{salv-} root four times, \latin{semper} twice, \latin{perpétu-} twice ---
so \textbf{the sung and read elements state the problem in the language of
service, trust and anxiety while the priest's own prayers state it in the
language of expiation and rescue.} The Alleluia's \latin{salutári} is the
single word tying the day's one purely jubilant chant into that thread.}
\pfield{Fruit}{The Mass has two registers of address, and a reader who follows
only the readings hears one of them. It also explains why the Alleluia,
otherwise the least integrated element of the formulary, is not simply
decorative.}
\pfield{What the element-by-element reading misses}{Read in order, the orations
are three short prayers in a conventional idiom. Read together and against the
chants, they are the only place in the Mass where the problem is named as
expiation.}
\pfield{Strongest limit}{\textbf{Severe, and it cannot be closed here.} This
is the ordinary vocabulary of Roman collects, so the saturation may be a fact
about the genre and not about this Mass --- and \textbf{the counter-test that
would settle it cannot be run}: a lexical concordance over the orations of the
whole Missal would have to run against an optical layer that demonstrably
undercounts, so its result would not be evidence. What can be said is bounded to
the fifteen formularies collated for this collection: the substantive
\latin{propitiátio} appears in the appointed Latin of \textbf{no other} of them.
That counter-test holds narrowly: its one apparent counter-example is English
audit prose in another leaf's record and not appointed text. \textbf{The wider \latin{propiti-} family is a different matter and the
count over it is not restated here}: the figure this guide carried for it comes
from an earlier sweep that no later one has repeated, so \textbf{the claim is
about the substantive and collapses if restated about the propitiation
word-family generally.} \textbf{Four words shared with the very next Sunday's
Collect --- \latin{Ecclésiam tuam}, \latin{quia sine te}, \latin{semper}, and
the two-verb doublet whose second verb is \latin{munire} --- are not
distinctive of this formulary either, so a claim of saturation must be made
about the words that are not shared.} \textbf{And the four-fold \latin{salut-} count is one root in four
senses} --- of a victim, of God, of things that save and of a state. Finally,
\textbf{this proposal is invisible in the guide's own English witness}: the
registered 1861 rendering drops \latin{propitiátio} at the Collect, turns it into
a verb at the Secret, and drops both \latin{semper} and the substantive at the
Postcommunion, so \textbf{every measurement supporting it is conducted on the
Latin}. That is how the 1861 translator rendered the three prayers, read on the
printed page at pp.~425 and 427, and not a doubt about what he wrote.}
\end{proposal}

\begin{proposal}{P4. The hope-formula lies past the cut, and is handed on at another Mass}
\pfield{Anchors}{Introit \cue{Int.}, Ps.\ 83:10--11a with vv.~2--3a; Gradual
\cue{Grad.}, Ps.\ 117:8--9; Offertory \cue{Off.}, Ps.\ 33:8--9a}
\pfield{Mechanism}{The formula \latin{beatus \dots\ qui sperat} stands in
exactly two verses of the Clementine Vulgate, and they are the two psalms this
Mass takes its Offertory and its Introit from: \latin{beatus vir qui sperat in
eo} at Ps.\ 33:9b and \latin{beatus homo qui sperat in te} at Ps.\ 83:13, the
two differing in what is hoped in, so that no one quoted form stands at both.
It lies \textbf{past the Introit antiphon's cut at v.~11a and past its psalm
verse's cut at v.~3a}; it lies \textbf{one clause past the Offertory's cut at
v.~9a}, in the second half of the very verse the antiphon ends on; and
\textbf{the Gradual alone is uncut and sings the hope vocabulary outright}, the
\latin{sper-} and \latin{confid-} stems standing twice each in the Gradual block
and not at all in the Introit. Three appointed excerpts stop short of the
formula and a fourth sings the vocabulary it is made of. \textbf{The Introit
belongs because the formula lies beyond both its boundaries, not because it uses
hope vocabulary --- which it does not.}}
\pfield{Fruit}{The formulary's chief coherence is a shared editorial boundary,
and it is verifiable rather than thematic. It also gives the Gradual, the day's
plainest maxim, a structural role it does not obviously have.}
\pfield{What the element-by-element reading misses}{An extent read one element
at a time is a boundary and nothing more: the Introit simply stops, the
Offertory simply stops. \textbf{Read across three chants the stopping is a
pattern, and the term the pattern stops before is appointed whole at another
Mass in the same book} --- the Eighth Sunday after Pentecost's Communion prints
Ps.\ 33:9 entire, \latin{Gustate et videte, quoniam suavis est Dominus: beatus
vir, qui sperat in eo}. Across the twenty-one collated text records covering
fifteen 1962 identities, \latin{sperat in eo} stands in the appointed Latin of
the Eighth Sunday alone and \latin{gustate} in the Eighth and the Fourteenth
only. So the beatitude of the one who hopes is not absent from the repertory; it
is assigned to the other Mass.}
\pfield{Strongest limit}{\textbf{The same datum has already been used, in this
collection, as a disconfirming condition rather than as evidence.} The Eighth
Sunday's own guide treats a pericope stopping one clause short, with the missing
clause appointed at another Mass in the same book, as showing that the withheld
matter ``is therefore not withheld from the liturgical year at all --- only from
this Sunday, and only because a lection had to end somewhere,'' which collapses
such a reading into a coincidence of scissors; and it records that
\latin{gustate et videte} is ``the commonest of Communion antiphons.''
\textbf{Both readings of one kind of evidence cannot be right, and this proposal
stands only so far as it answers that objection.} Beyond it: \textbf{the
antiphons' extents are inherited chant tradition and nothing establishes that
anything was cut in order to suppress the formula}; psalms of supplication
routinely close on exactly such promises, so the pattern may be a property of
the psalter rather than of the selection; duller explanations are available,
metrical length and the ordinary bounds of an antiphon among them; and
\textbf{the test that would settle it --- measuring how often an arbitrary chant
extent in this repertory stops short of a comparable promise --- cannot be run
on the sources held here, which include no chant repertory at all.} Secondly,
\latin{confidere} is a trust word rather than a hope word strictly: a stricter
count of \latin{spes} vocabulary alone finds \latin{speráre} twice in the
Gradual and nothing else. Thirdly, the Eighth Sunday's Communion prints
Ps.\ 33:9 whole with the hope clause in the Pustet Ratisbon 1862 and the Venice
1570 text layers, both searched rather than inspected and \textbf{neither
collated against a page image.} Fourthly, Ambrose quotes the fuller verse
\emph{with} the hope clause in the fourth century, which sharpens the
observation about the Roman cut but says nothing about why it was made.}
\end{proposal}

\begin{proposal}{P5. Two chants, two divisions: the Fathers disagree at the Offertory and at the Gradual, and the formulary resolves neither}
\pfield{Anchors}{Offertory \cue{Off.}, Ps.\ 33:8--9a, \latin{Immíttet Angelus
Dómini \dots\ gustáte et vidéte}; Gradual \cue{Grad.}, Ps.\ 117:8--9,
\latin{quam speráre in princípibus}}
\pfield{Mechanism}{The Fathers divide at each of the two chants, and the two
divisions are \textbf{independent} of each other; the Mass prints both without
adjudicating. At the
\textbf{Offertory} the angel is read four ways --- Augustine's Christ, the Angel
of great counsel; Cassiodorus's covert working and moral sense, on the very
lemma Augustine denounced; Arnobius's angel of the sepulchre; Godfrey's gift of
grace --- while the same verse is attested by the \work{Mystagogical Catecheses}
and by Ambrose as the Church's \emph{communion} invitation, and the 1962 book
assigns it to the \emph{offertory}. At the \textbf{Gradual} Augustine escalates
the excluded second term to \textbf{good angels} --- ``for angels also are called
princes,'' in the English of the \work{Nicene and Post-Nicene Fathers}, first
series --- making the exclusion absolute rather than moral, against
Theodoret's temporary earthly rulers and Godfrey's interior and exterior senses.
\textbf{The angel who is the Mass's rescuer at one chant is the Mass's excluded
object of trust at the other, and the Fathers reached each position
separately.}}
\pfield{Fruit}{A cross-chant tension the Fathers themselves hold and the
formulary's own arrangement stages; and a way to treat
the Offertory's angel that neither flattens the christological
identification nor asserts it.}
\pfield{What the element-by-element reading misses}{Each division is reported
under its own chant as a local disagreement. Set side by side, they are the
same word doing opposite work at two moments of one Mass.}
\pfield{Strongest limit}{\textbf{Both divisions are deeper than when this was
drafted, and one of them may be shallower than it looks.} \textbf{Arnobius's
angel of the sepulchre depends on repointing \latin{timéntium eum} as those
fearing \emph{for} his body, which is not what the missal's Latin says.}
\textbf{And the Gradual's angelic reading is one position with three
carriers and not three witnesses.} Augustine holds it at his own exposition of
Ps.\ 117, with Daniel's Michael as warrant and with the principle that
reconciles the two chants: no one is to be trusted in, ``for no one is good,
save God alone; and when a man or an angel appear to aid us \dots\ He does it
through them, who made them good after their measure.'' Peter Lombard puts the
reading under Augustine's name with the same warrant, so he transmits him
correctly; Prosper's exposition is largely an epitome of Augustine, and
Prosper's reason is Lombard's reason. \textbf{But that reconciliation is
conditional on who the Offertory's angel is, and for Augustine the tension never
arises at all, because for him that \latin{Angelus} is Christ. If the angel is a
creature --- as Cassiodorus, Godfrey and Arnobius have it --- the Mass does trust
at one chant a rescuer it excludes from trust at the other, and the instrumental
principle softens the tension without removing it.} Cassiodorus, Godfrey and Arnobius are read here in Latin, in
transcriptions of \work{Patrologia Latina} 70, 174 and 53, with no page image
collated and no critical edition beside them; three witnesses of this proposal
stand in translation --- Augustine, whose \work{Enarratio} on Ps.\ 117 is
quoted in the English of the \work{Nicene and Post-Nicene Fathers}, first
series, with \work{Patrologia Latina} 36--37 uncollated, and the
\work{Mystagogical Catecheses} and Ambrose, read in the second series of the
same collection, with \work{Patrologia Graeca} 33 and \work{Patrologia Latina}
16 uncollated. Bruno the Carthusian's
question is a fifth thing and is grammatical rather than doctrinal.
\textbf{The cross-relation between the Gradual's absolute exclusion and the
Gospel's \latin{Nemo potest duóbus dóminis servíre} is this guide's observation
and not Augustine's: he does not cite Mt.\ 6:24 at that locus}, and the Gospel is
accordingly not an anchor here. The Offertory half carries three further limits:
Augustine's text differs from the missal's at the following verb, \latin{éruet}
against \latin{erípiet}, which is evidence that the missal's text is \textbf{not}
simply his; the missal's capital \latin{Angelus} is orthography and \textbf{does
not encode} his identification; and the attribution of the
\work{Mystagogical Catecheses} to Cyril of Jerusalem is disputed in modern
scholarship and that question was not swept.}
\end{proposal}

\begin{proposal}{P6. The one verb the Mass appoints twice, with its object reversed}
\pfield{Anchors}{Introit \cue{Int.}, the psalm verse, \latin{concupíscit, et
déficit ánima mea in átria Dómini}; Epistle \cue{Ep.}, \latin{caro enim
concupíscit advérsus spíritum} and \latin{carnem suam crucifixérunt cum vítiis
et concupiscéntiis}}
\pfield{Mechanism}{\latin{Concupiscere} is appointed three times in two
elements with its object reversed. In the Introit's psalm verse the soul longs
\emph{toward} the courts of the Lord and faints for them; in the Epistle the
flesh desires \emph{against} the spirit, and what is crucified is the flesh
\emph{with its desires}. \textbf{One verb states both the movement the Mass
asks for and the movement it asks to be delivered from}, and the difference
between them is the direction of the longing rather than its intensity.}
\pfield{Fruit}{Desire is not the enemy in this formulary; its object is. That
gives the Epistle's mortification a positive term inside the same Mass and
keeps the crucifixion of the flesh from being read as the extinction of
longing.}
\pfield{What the element-by-element reading misses}{The Introit's occurrence
sits in the psalm verse, which an element-by-element reading treats as the
antiphon's setting rather than as appointed vocabulary; the Epistle's two are
read as Pauline anthropology. Nothing brings the three together.}
\pfield{Strongest limit}{\textbf{Two limits, both real, and a third that
forbids the reinforcement one would want.} The two senses are the ordinary
senses of an ordinary Latin verb, \textbf{so the opposition may be a
coincidence of the Vulgate's vocabulary before it is a compositional choice,
and no compiler's intention is claimed.} The Introit's occurrence is in the
\textbf{psalm verse} and not the antiphon, so it is liable to the objection that
a psalm verse is not appointed in the same sense the antiphon is. \textbf{And
the second lexical leg that would strengthen it is not used here:}
\latin{ánima} stands in the same Introit clause, three words from
\latin{concupíscit}, and again in the Gospel with the valuation reversed ---
\textbf{but \latin{anima} is among the commonest nouns of both psalter and
Gospels, and Mt.\ 6:25's \latin{ánimæ vestræ} is very likely ``your life''
rather than ``your soul,'' in which case the Introit's longing soul and the
Gospel's anxious life are two different subjects and the reversal is a trick of
Latin vocabulary.} \latin{Concupíscit} is also one of only five substantive
forms a mechanical comparison of the ten appointed blocks returns as shared
between different texts, so it carries no more weight than the other four.}
\end{proposal}

\clearpage
